\section{The evidence-binding gap in AI reasoning systems}

AI systems that generate natural language outputs about scientific or physical domains make claims
that require empirical grounding. When a language model responds to a question about orca encounter
likelihood at Haro Strait, it may correctly name the relevant salmon corridor, acoustic monitoring
network, and behavioral studies --- but without binding any of these claims to their generating
observations, prior analyses, or data artifacts. The response is fluent and plausible; it is not
evidentially accountable.

This is the evidence-binding gap: the structural disconnection between a system's stated claims and
the observations, methods, and experiments that would make those claims verifiable. The gap is not a
matter of factual accuracy per se --- a claim may be broadly correct while still carrying no bound
evidence. It is a matter of verifiability: a judge, steward, or downstream system cannot determine
whether the claim was produced by accessing ground-truth data or by reproducing a plausible-sounding
pattern from training.

Prior work on RAG hallucination reduction~\cite{arxiv230915217,arxiv230514627,arxiv240115884,arxiv231109533}
establishes that retrieval-augmented generation reduces factual error rates and that citation generation
can be automated. However, this literature focuses on whether RAG improves accuracy, not on whether
the improvement comes from evidence binding --- from the system's outputs being traceable to their
generating context --- versus improved parametric knowledge activation.

The geospatial grounding baseline makes this distinction concrete. Google Maps grounding achieves high
citation density for place-of-interest queries. For marine science evidence queries --- the same domain,
but requiring scientific rather than place-level grounding --- Maps grounding resolves 0 of 25 citations
to scientific or dataset sources. NOAA fisheries data, the Center for Whale Research census, Chinook
salmon migration corridors, and acoustic hydrophone recordings are all generated as claims but remain
entirely uncited. The system knows about these phenomena from training; it cannot bind its claims to
verifiable data artifacts.

$\Runcited$ --- the fraction of scientific sentences in a system's output with no bound
citation --- is the measurement primitive that makes this gap quantitative, reproducible, and
architecture-diagnostic.
